\documentclass[11pt]{article}
\usepackage[T1]{fontenc}
\usepackage{lmodern}
\usepackage[margin=0.85in]{geometry}
\usepackage{amsmath,amssymb,amsthm}
\usepackage[hidelinks]{hyperref}
\newtheorem{theorem}{Theorem}
\newtheorem{corollary}[theorem]{Corollary}
\newcommand{\F}{\mathbb F_2}
\newcommand{\supp}{\operatorname{supp}}
\newcommand{\im}{\operatorname{im}}

\title{Boolean polynomial obstructions to binary expansion\newline
of balanced Cayley complexes}
\author{Kaiwen Jin\\\small\texttt{yui.ui.siki@gmail.com}}
\date{October 14, 2026}

\begin{document}
\maketitle

\begin{abstract}
The proposed coboundary expansion of random balanced Cayley complexes
with binary coefficients and logarithmically many group-product layers
fails in every fixed dimension $k\ge2$. For $G=(\F)^d$ and every
$A\subseteq G$, we prove
\[
\dim_{\F}H^{k-1}(Y_{A,k};\F)
\ge \sum_{j=0}^{\min(k,d)}\binom dj-|A|.
\]
The proof maps a Boolean polynomial vanishing on $A$ to a cohomology class
by writing the polynomial of a total group sum as a simplicial coboundary
on the complete join. Consequently, for any fixed positive constant $C$,
every subset of size $\lceil C\log D(G)\rceil$ has zero binary coboundary
expansion for all sufficiently large $d$, where $D(G)$ is the sum of the
degrees of its irreducible complex representations.
\end{abstract}

\section{The complex and the obstruction}

Coboundary expansion measures the number of violated parity constraints
relative to the distance of a binary cochain from the coboundary space.
Meshulam introduced balanced Cayley complexes and proved a real-coefficient
spectral-gap estimate when the number of allowed group products is
logarithmic in a representation parameter \cite{meshulam}.
Conjecture~8.1 of that paper asks for a binary coboundary-expansion analogue,
uniformly over finite groups whose orders tend to infinity.
We give a deterministic obstruction on elementary abelian $2$-groups.

Fix integers $d\ge1$ and $k\ge2$, and write $G=(\F)^d$ additively.
Let $V_i=\{i\}\times G$ for $1\le i\le k+1$, and put
$X=V_1*\cdots*V_{k+1}$. A face of this join contains at most one vertex
from each part. For $A\subseteq G$, let $Y_{A,k}$ contain the complete
$(k-1)$-skeleton of $X$ and exactly the $k$-faces with labels
$(g_1,\ldots,g_{k+1})$ satisfying
\[
g_1+\cdots+g_{k+1}\in A.
\]
Addition here is the ordered group product of the original construction.

All cochains below have coefficients in $\F$. We use the increasing order
of vertex parts to identify a face with its tuple of labels.
The simplicial coboundary $\delta_r:C^r(Y;\F)\to C^{r+1}(Y;\F)$ adds the
values on the codimension-one faces. Write
\[
B^r(Y)=\im\delta_{r-1},\qquad
Z^r(Y)=\ker\delta_r,\qquad H^r(Y;\F)=Z^r(Y)/B^r(Y)
\]
for $r\ge1$. For a cochain $\phi$, let $|\phi|=|\supp\phi|$ be its
unnormalized support count. The expansion constant under consideration is
\begin{equation}\label{eq:expansion}
h_{k-1}(Y)=\min_{\phi\notin B^{k-1}(Y)}
\frac{|\delta_{k-1}\phi|}{\min_{b\in B^{k-1}(Y)}|\phi+b|}.
\end{equation}
A nonzero cohomology class has a representative with zero numerator and
strictly positive denominator in \eqref{eq:expansion}. Thus nonvanishing
of $H^{k-1}(Y;\F)$ forces $h_{k-1}(Y)=0$.

\begin{theorem}\label{thm:obstruction}
For every $d\ge1$, $k\ge2$ and $A\subseteq(\F)^d$, put
\[
N_k(d)=\sum_{j=0}^{\min(k,d)}\binom dj.
\]
Then
\begin{equation}\label{eq:dimension}
\dim_{\F}H^{k-1}(Y_{A,k};\F)\ge N_k(d)-|A|.
\end{equation}
In particular, $|A|<N_k(d)$ implies $h_{k-1}(Y_{A,k})=0$.
\end{theorem}

For a finite group $G$, write
$D(G)=\sum_{\rho\in\widehat G}\dim\rho$, where $\widehat G$ denotes its
irreducible complex representations, as in \cite{meshulam}.
For the groups considered here, $D((\F)^d)=2^d$, as proved below,
so $\log D(G)=d\log2$.

The number $N_k(d)$ is the dimension of the Boolean polynomial functions
of degree at most $k$ on $(\F)^d$. Vanishing on $A$ imposes at most $|A|$
linear constraints on these functions. The central point is that the
remaining functions give distinct cohomology classes in the specified
complex. Since $N_k(d)$ is at least quadratic in $d$ for $k\ge2$, this
already rules out logarithmic sampling.

\begin{corollary}\label{cor:logarithmic}
Fix $k\ge2$ and $C,\varepsilon>0$. For $G_d=(\F)^d$, let $A$ be uniform
among the subsets of size $m_d=\lceil C\log D(G_d)\rceil$.
For all sufficiently large $d$, this sample size is admissible and
\[
\Pr\{h_{k-1}(Y_{A,k})\ge\varepsilon\}=0.
\]
Consequently, the constants requested in the binary expansion conjecture
do not exist for any fixed $k\ge2$.
\end{corollary}

\section{Proofs}

\begin{proof}[Proof of Theorem~\ref{thm:obstruction}]
Let $\mathcal P_k(d)$ be the space of Boolean polynomial functions
\[
q(t)=\sum_{\substack{S\subseteq[d]\\|S|\le k}}c_S\prod_{r\in S}t_r,
\qquad c_S\in\F,
\]
where $[d]=\{1,\ldots,d\}$ and the empty product is one.
These monomials are independent as functions. Indeed, evaluation at the
vector $1_T$ supported on $T$ gives $q(1_T)=\sum_{S\subseteq T}c_S$.
For a zero function, the empty set first gives $c_\varnothing=0$, and
induction on $|T|\le k$ then gives every $c_T=0$.
Therefore $\dim\mathcal P_k(d)=N_k(d)$.

Let $K_A$ be the subspace of polynomials vanishing on $A$.
The evaluation map $\mathcal P_k(d)\to\F^A$ has rank at most $|A|$, so
\begin{equation}\label{eq:kernel}
\dim K_A\ge N_k(d)-|A|.
\end{equation}
We construct a linear injection from $K_A$ into the cohomology in
\eqref{eq:dimension}.

For each $q\in\mathcal P_k(d)$, expand $q(g_1+\cdots+g_{k+1})$.
For a monomial indexed by $S$, the expansion is
\begin{equation}\label{eq:monomial}
\prod_{r\in S}\left(\sum_{j=1}^{k+1}(g_j)_r\right)
=\sum_{\alpha:S\to[k+1]}\prod_{r\in S}(g_{\alpha(r)})_r.
\end{equation}
Each term uses at most $|S|\le k$ of the $k+1$ vertex parts.
Assign it to the smallest part index $i$ outside $\im\alpha$;
for the empty monomial assign it to $i=1$.
Grouping the assigned terms, with their coefficients $c_S$, yields functions
$f_i^q$ independent of $g_i$ such that
\begin{equation}\label{eq:decomposition}
q(g_1+\cdots+g_{k+1})
=\sum_{i=1}^{k+1}f_i^q(g_1,\ldots,\widehat{g_i},\ldots,g_{k+1}).
\end{equation}
The assignment is fixed independently of $q$, so the functions depend
linearly on its coefficients.

Every $(k-1)$-face of $X$ misses exactly one vertex part. Define $Tq$ on
a face missing part $i$ to have value $f_i^q$ at that face's labels.
Equation~\eqref{eq:decomposition} is precisely the identity
\begin{equation}\label{eq:coboundary}
\delta_X(Tq)(g_1,\ldots,g_{k+1})=q(g_1+\cdots+g_{k+1}).
\end{equation}
If $q\in K_A$, the right-hand side vanishes on every top face of
$Y_{A,k}$. Hence $Tq$ is a cocycle there, and $q\mapsto[Tq]$ is a linear
map into $H^{k-1}(Y_{A,k};\F)$.

To prove injectivity, suppose $Tq=\delta_{k-2}\psi$ in $Y_{A,k}$.
The complexes $Y_{A,k}$ and $X$ have identical faces and incidences through
dimension $k-1$, so the same equality holds in $X$.
Applying its next coboundary gives $\delta_X(Tq)=0$:
in $\delta^2$, every codimension-two subface occurs twice, and $2=0$ in
$\F$. Equation~\eqref{eq:coboundary}, evaluated at $(t,0,\ldots,0)$,
now gives $q(t)=0$ for every $t\in G$. This is a valid top face of $X$
because its vertex parts are disjoint, even when some labels coincide.
Thus $q=0$, proving injectivity. Together with \eqref{eq:kernel}, it proves
\eqref{eq:dimension}.

If $|A|<N_k(d)$, choose a nonzero $q\in K_A$.
Then $Tq$ is a cocycle outside $B^{k-1}(Y_{A,k})$.
Each $Tq+b$, with $b$ in that coboundary space, is nonzero, so its support
has at least one face. The denominator in \eqref{eq:expansion} is thus a
finite positive minimum, while the numerator is zero. All ratios are
nonnegative, and the asserted zero expansion follows.
\end{proof}

\begin{proof}[Proof of Corollary~\ref{cor:logarithmic}]
Every irreducible complex representation of $G_d$ is one-dimensional.
Indeed, its coordinate generators act by commuting involutions.
For an involution $M$, the projectors $(I+M)/2$ and $(I-M)/2$ split into
its two eigenspaces. The other generators preserve both spaces, and
successive splitting gives simultaneous eigenspaces. Every line in a
nonzero simultaneous eigenspace is invariant, so irreducibility forces
dimension one. The $2^d$ independent choices of signs on the coordinate
generators give exactly the characters
$\chi_u(t)=(-1)^{u\cdot t}$ for $u\in(\F)^d$. Consequently $D(G_d)=2^d$.

Choose $d_0\ge2$ so that $(d_0+1)/2>C\log2$.
For every $d\ge d_0$, using natural logarithms,
\[
\begin{aligned}
m_d=\lceil C d\log2\rceil
&\le C d\log2+1
<1+\frac{d(d+1)}2\\
&=1+d+\binom d2\le N_k(d)\le2^d.
\end{aligned}
\]
The sample size is therefore admissible. Theorem~\ref{thm:obstruction}
gives $h_{k-1}(Y_{A,k})=0$ for every subset of that size, making the
displayed probability zero for every $\varepsilon>0$.
Since $|G_d|=2^d\to\infty$, this sequence violates the required guarantee
for every proposed pair of constants and each fixed $k\ge2$.
A different fixed base of the logarithm only rescales $C$.
\end{proof}

\paragraph{AI-use statement.}
OpenAI Codex assisted with mathematical derivation, proof checking and
drafting of this manuscript. The mathematical claim and the complete proof
are presented here for independent review.

\begin{thebibliography}{1}
\bibitem{meshulam}
R.~Meshulam, \emph{Random Balanced Cayley Complexes},
arXiv:2211.02085v1 (2022), Sections~1--2 and Conjecture~8.1.
\url{https://arxiv.org/abs/2211.02085v1}.
\end{thebibliography}

\end{document}
